\documentclass[11,5pt, a4paper]{article}
\usepackage[a4paper, total={6in, 9.5in}]{geometry}
\usepackage{float}
\usepackage{graphicx} % Required for inserting images
\usepackage{natbib} % bibliografia
\usepackage{url}
\usepackage{booktabs} % tabella
\usepackage{amsmath}

\usepackage[italian]{babel}
\usepackage[utf8]{inputenc} % Per caratteri accentati
\usepackage[T1]{fontenc}    % Per font corretti in output

\title{\Large Analisi della sottoscrizione di un conto deposito \\[0.5em] \large Introduzione alla Data Science e al Pensiero Computazionale}


\author{\large Gruppo 14 \\[0.33em] Arianna Bertoldo e Veronica Contarin}
\date{15 luglio 2026}

\begin{document}

\maketitle

\tableofcontents %COME SI METTE IN ITALIANO :(

\section{Introduzione}
% che attacco tremendo
% il telemarketing spacca
Il telemarketing è una pratica diffusa in numerosi settori, tra cui quello bancario e finanziario. Consiste nel contattare direttamente potenziali clienti per promuovere prodotti o servizi. 
Sebbene sia un metodo rapido e potenzialmente efficace, comporta costi organizzartivi ed economici significativi, con il rischio di disperdere risorse su contatti infruttuosi. La sua efficacia dipende infatti dalla capacità di individuare a monte i clienti probabilmente più propensi all'adesione o all'acquisto e contattarli nel momento più adeguato.

A questo proposito, la crescente disponibilità di dati consente e incoraggia le aziende a ricorrere a un approccio \textit{data-driven} per ottimizzare le loro attività.
Le informazioni raccolte nelle campagna possono orientare l'organizzazione delle campagne successive, selezionando accuratamente i destinatari e personalizzando le comunicazioni. L'analisi dei comportamenti e delle caratteristiche della clientela può così contribuire a ridurre i contatti infruttuosi e a rendere il telemarketing più mirato ed efficiente \citep{ROSARIO2023}.
\\

In questa relazione, presentiamo un'analisi dei dati \citep{ariannab02}
che abbiamo condotto sulla base dei risultati di una campagna di telemarketing bancaria mirata alla sottoscrizione di conti deposito vincolati.  L'obiettivo ultimo era quello di predire quali clienti avrebbero sottoscritto il prodotto proposto, costruendo modelli predittivi per identificare i clienti che presentano una maggiore probabilità di sottoscrivere il conto deposito. 


\section{Dataset}

Il dataset esaminato è \textit{Bank\_Marketing.csv} \citep{BankMarketing}.
Si tratta di un insieme di dati raccolti da una banca portoghese tra maggio 2008 e novembre 2010, nel corso di una campagna marketing per la sottoscrizione di un conto deposito vincolato \citep{bank_marketing_222}.

Il dataset si compone di 45211 osservazioni, per le quali sono raccolte informazioni relative a 17 variabili, di cui 16 variabili esplicative e una variabile target (\textit{y}). 
Le variabili esplicative possono essere suddivise in quattro gruppi. Il primo comprende le caratteristiche personali e socio-economiche del cliente, quali età, occupazione, stato civile e livello di istruzione. Il secondo riguarda la sua situazione finanziaria, attraverso il saldo medio annuo, la presenza di prestiti personali o immobiliari e l'eventuale inadempienza creditizia. Il terzo gruppo descrive la campagna corrente, includendo il mezzo di contatto, il giorno e il mese dell'ultima chiamata, la sua durata e il numero complessivo di contatti effettuati. Infine, alcune variabili riportano informazioni sulla precedente campagna, come il numero di contatti, il tempo trascorso dall'ultima chiamata e il relativo esito.

% sposto tabella


\begin{table}[h]
\centering
\caption{Classificazione delle variabili per area tematica}
\begin{tabular}{@{}ll@{}}
\toprule
\textbf{Gruppo} & \textbf{Variabili} \\ \midrule
Socio-demografiche & age, job, marital, education \\
Finanziarie & balance, default, loan, housing \\
Campagna corrente & contact, day, month, duration, campaign \\
Campagna precedente & pdays, previous, poutcome \\
Target & y \\ \bottomrule
\end{tabular}
\end{table}

Il dataset è ben strutturato, ma in alcune categorie presenta numerosi valori mancanti, categorizzati come \textit{unknown}. Per la variabile \textit{contact}, il 29\% dei valori sono assenti, poiché la raccolta di tale dato è iniziata dopo due mesi dall'avvio della campagna. Ulteriormente, anche \textit{job} e \textit{education} presentano delle basse quote (inferiori al 5\%) di valori mancanti. Infine, per \textit{poutcome}, quasi l'82\% delle osservazioni sono identificate come sconosciute; tuttavia, ciò non è un segnale di bassa qualità dei dati, bensì è legato alla natura della variabile: le osservazioni mancanti sono per lo più riconducibili a clienti che non sono stati contattati nella precedente campagna, e di conseguenza l'esito non è identificabile.

In seguito è stato notato che, in mancanza di un identificativo univoco per i singoli clienti contattati, c'è la possibilità che talune persone siano state incluse nel dataset più volte, a seguito di ripetute chiamate. In questo caso, le osservazioni non sono mai esattamente corrispondenti, se non per le caratteristiche socio-demografiche e finanziarie.

\section{Metodologia}
Il lavoro di analisi dei dati è stato organizzato in tre fasi. Nella prima è stata realizzata un'analisi descrittiva del dataset, mirata a comprenderne la struttura, la qualità e le distribuzioni delle variabili. Nella seconda fase sono state esaminate le relazioni tra le variabili esplicative e la variabile target, con il supporto di visualizzazioni grafiche. Infine, nella terza fase, i dati sono stati preparati per la classificazione e sono stati confrontati diversi algoritmi di apprendimento supervisionato.

\subsection{Fase 1: Analisi descrittiva ed esplorativa}
\label{ss: Fase 1}
La prima fase ha riguardato l'analisi della struttura del dataset, dei tipi di variabile, della presenza di informazioni mancanti, e di alcune relazioni tra le variabili. 

Per le variabili numeriche sono state calcolate statistiche descrittive, come misure di centralità e di dispersione. Per le variabili categoriche sono state invece considerate le frequenze assolute e relative.

Tramite l'analisi, sono emersi dei potenziali limiti del dataset, legati prevalentemente a valori mancanti (\textit{unknown}) e allo sbilanciamento della variabile target.

Ulteriormente, nella prima fase, sono state avanzate domande e ipotesi relative alle caratteristiche legate alla sottoscrizione. Queste relazioni sono state esplorate tramite confronti tra categorie e gruppi di clienti, e coefficienti di correlazione. In particolare, sono state esaminate le differenze nella propensione alla sottoscrizione secondo caratteristiche socio-demografiche e finanziarie, al numero di contatti effettuati, all'esito della precedente campagna e alla durata delle chiamate.

\subsection{Fase 2: Analisi delle relazioni tra le variabili}
\label{ss:Fase2}
Nella seconda fase, l'analisi si è concentrata sulle relazioni tra le variabili esplicative e la variabile target, ossia l'esito della campagna. 
L'obiettivo era individuare quali caratteristiche dei clienti e quali aspetti dei contatti fossero associati a un esito positivo della campagna. Le analisi svolte in questa sezione si sono innestate sui risultati e riflessioni della \ref{ss: Fase 1}.

Sono stati confrontati i tassi di adesione tra diverse categorie lavorative, situazione finanziaria, fasce d'età e mezzo di contatto. Inoltre, sono state esaminate le relazioni tra l'esito della campagna e le variabili relative alla precedente campagna, tra cui il numero di chiamate già effettuate, il numero di contatti ed esito della campagna precedente. Un'ulteriore aspetto esplorato è stata la durata delle chiamate, con la sua distribuzione e andamento medio, comparandole in base all'esito della campagna. 
A questo proposito, per porre enfasi sulla successione temporale, la variabile \textit{month} è stata trasformata in un indice progressivo. In tal modo, mesi con lo stesso nome ma appartenenti ad anni diversi non sono stati ridotti alla stessa categoria, così facilitando la rappresentazione dell'evoluzione dell'organizzazione della campagna nel tempo.

Infine, si osserva che le categorie più positivamente associate alla variabile target sono sottorappresentate nel dataset, limitando la portata dei risultati.

\subsection{Fase 3: Modelli di classificazione}
\label{ss: Fase 3}
In questa fase si sono confrontati tre modelli di classificazione supervisionata: regressione logistica, k-Nearest Neighbours e Random Forest.

Il primo modello, la regressione logistica, stima la probabilit\`{a} che un'osservazione appartenga alla classe positiva, esaminando la relazione tra variabili predittive e variabile target \citep{REGLOG_IBM}. Per normalizzare le variabili sulle quali \`{e} stato allentato il modello, \`{e} stata utilizzata la seguente formula per la standardizzazione:

\begin{equation*}
    Z = \frac{X-\mu}{\sigma}
\end{equation*}

Il secondo modello realizzato è k-Nearest Neighbours. Questo modello assegna una classe a ciascuna osservazione in base alle classi prevalenti tra le $k$ osservazioni più vicine all'interno del dataset. Dato che la sua performance dipende dal numero di vicini presi in considerazione, sono stati confrontati vari valori di $k$ e per ciascuno di questi sono stati valutati gli effetti su f1 score per la classe positiva e accuracy complessiva. Sulla base dei risultati ottenuti, il modello finale è stato addestrato utilizzando $k=3$.

Infine, è stato considerato il modello Random Forest. Tale modello è composto da un insieme di alberi decisionali, ciascuno addestrato su differenti campioni e sottoinsiemi di variabili. La combinazione di più alberi rende possibile arrivare a previsioni più robuste rispetto a quelle prodotte da un singolo albero decisionale, grazie all'ensemble learning \citep{RF_IBM}.
Per la Random Forest sono stati confrontati diverse quantità di stimatori e soglie di classificazione. Alla fine, si è optato per un modello con 100 alberi e una soglia di 0.6. Inoltre, sono stati aggiustati alcuni parametri, come la profondità, e si è optato per dare un peso bilanciato alle osservazioni (\textit{weight=balanced\_subsample}), in modo da assegnare un peso maggiore alla classe positiva che è sottorappresentata nel dataset.

\section{Dati e Interpretazione}
 %risultati per le varie fasi o complessivi
 %* grafici (2+)
 %* tabella (1+) es. confronto modelli
\subsection{Fase 1 e 2: Domande di ricerca e Analisi}
Dalle analisi esplorative è emerso fin da subito come la classe target (i clienti iscritti ad un conto deposito presso la banca) sia in forte minoranza rispetto ai clienti non iscritti all'interno del campione: al momento della raccolta dei dati, solo l’11\% del totale ha aderito ad un conto deposito. 

\subsubsection{Classi lavorative, giacenza media, livello di istruzione}

Ci si è domandati se l'appartenenza a determinate classi lavorative potesse portare ad un maggiore desiderio di iscriversi ad un conto deposito vincolato.
È emerso che gli studenti e i pensionati sono state le classi lavorative più iscritte, rispettivamente il 28,7\% e il 22,8\% di loro. Si è verificato che giovani e anziani avessero effettivamente tassi di iscrizione maggiori rispetto alla fascia intermedia e così è risultato: il 33,4\% dei clienti over 60 è iscritto ad un conto deposito, il 17,6\% degli under 30, e solo il 10\% di chi ha fra i 30 e i 59 anni è iscritto.

\begin{figure}[H]
\centering
\includegraphics[width=0.5\textwidth]{GRAFICI PER LATEX/subscription-jobs.png}
\caption{Distribuzione del tasso di iscrizione a conto deposito per classi lavorative}
\label{fig:job}
\end{figure}

Gli operai sono stati il gruppo professionale meno iscritto, con appena il 7,3\% di adesioni, seguiti da imprenditori, casalinghi e lavoratori nei servizi, tutti con un tasso di iscrizione che non supera il 9\%.
Ipotizzando che la minore adesione potesse essere dovuta ad un fattore condiviso fra queste classi lavorative, sono state vagliate due ipotesi: \textit{è la disponibilità economica ad agire sulla propensione ad aderire ad un conto deposito o il livello di istruzione?}

Nessuna delle due è risultata in grado di spiegare il comportamento della variabile target (tenuto conto, cioè, che i più propensi ad aprire un conto deposito sono stati studenti, pensionati, disoccupati, manager e amministrativi).
È vero che le categorie con saldi medi annui più alti sono più iscritte (pensionati), ma non si tratta di un rapporto causa-effetto: gli studenti, ad esempio, sono fra una delle due classi più iscritte, e pur essendo i più risparmiatori (i loro conti sono i meno in negativo fra tutte le classi lavorative), %eppure 
non sono i più ricchi. Le categorie con saldi medi annui più bassi sono effettivamente meno interessate ad aprire un conto deposito (vedi operai), ma non c’è un trend preciso e identificabile.

Non notiamo una particolare relazione nemmeno fra livello di istruzione e tasso di iscrizione. I più istruiti a livello universitario sono risultati manager, lavoratori autonomi, imprenditori, mentre lavoratori nei servizi e tecnici sono le classi con le percentuali più alte di diplomati. Tuttavia questi gruppi non si sono distinti per una maggiore propensione verso l’adesione allo strumento finanziario propostogli dalla banca.
E in effetti i clienti con un livello di istruzione superiore alla scuola primaria risultano solo leggermente più propensi a sottoscrivere un conto deposito. In particolare, il 15\% di chi possiede un’istruzione universitaria o equiparabile è iscritto, contro l’8,6\% di chi possiede la sola licenza elementare.

\begin{figure}[H]
\centering
\includegraphics[width=0.6\textwidth]{GRAFICI PER LATEX/education-jobs.png}
\caption{Distribuzione del livello di istruzione per classi lavorative}
\label{fig:education-job}
\end{figure}

Concludiamo piuttosto, allora, che studenti e pensionati siano profili più in linea con lo strumento finanziario del conto deposito vincolato. Ad iscriversi potrebbe spingerli la ricerca di uno strumento semplice e affidabile per i propri risparmi, ma questo aspetto non è indagabile da una ricerca come quella svolta in questa sede.

\subsubsection{Campagna precedente}
Altri aspetti sono risultati più rilevanti. È stato studiato \textit{il rapporto tra la campagna precedente e quella corrente}.
I clienti che già hanno sottoscritto un conto deposito hanno tassi di successo molto più alti rispetto agli altri gruppi.
Il 23\% dei clienti già contattati nella campagna marketing precedente si sono iscritti ad un conto deposito, contro il 9\% dei clienti mai contattati.
Se la scorsa campagna aveva già avuto successo, la percentuale di iscrizioni in quella attuale arriva al 64,7\%. I clienti iscritti ma non convinti nella campagna passata sono invece il 12,6\%.
Inoltre, un alto numero di contatti nella storia del cliente è associato ad un tasso di iscrizione più alto.

\begin{figure}[H]
\centering
\includegraphics[width=0.6\textwidth]{GRAFICI PER LATEX/contact-frequency-poutcome-success.png}
\caption{Tassi di successo e insuccesso per numero di contatti ed esito della campagna precedente}
\label{fig:frequency-contact-poutcome}
\end{figure}

Clienti già contattati in passato, a maggior ragione se hanno già aderito ad un'offerta della banca, dimostrano una fiducia nei confronti della banca che si concretizza in tassi di adesione più alti alle sue campagne marketing.

Diversamente da quella precedente, nella campagna attuale, la probabilità di successo è massima per i clienti contattati una singola volta: 14,6\% per una chiamata fino a scemare al 4\% per più di undici chiamate.
La durata media di una chiamata è inferiore ai cinque minuti. La durata massima è stata di un'ora e venti.
Tolti i casi limite, le chiamate più lunghe sembrano aggirarsi attorno agli undici minuti e quelle standard sui tre.

\begin{figure}[H]
\centering
\includegraphics[width=1\textwidth]{GRAFICI PER LATEX/duration-success-rate.png}
\caption{Andamento tasso successo e insuccesso delle chiamate per durata}
\label{fig:duration-success}
\end{figure}

Emerge una chiara associazione tra durata delle chiamate e tasso di successo: le chiamate con hanno tassi di adesione più alti sono quelle con durate nettamente maggiori. Questo andamento è osservabile soprattutto nei primi tre quarti della campagna, ma persiste anche nella fase finale, dove l'andamento delle chiamate di successo si conforma a quello delle chiamate senza successo. 
Ipotizziamo che a quel punto il numero di clienti disponibili ad essere chiamati sia diminuito ulteriormente, portando ad un allineamento con la durata media delle chiamate per il campione (campione, ricordiamo, fortemente sbilanciato per la sovrabbondanza di non iscritti).

Abbiamo notato infatti che man mano che le chiamate si allungano queste sono associate a percentuali di sottoscrizioni più alte: a cinque minuti, almeno un cliente su dieci sottoscrive; a dieci minuti, sono diventati quattro;  quindici, sei clienti su dieci risultano iscritti.

\subsubsection{Periodo della campagna}

Il maggior numero di chiamate è concentrato nel primo anno, con picchi intorno ad alcuni intervalli. La media in questo primo periodo è di circa 3000 chiamate al mese. Nel secondo anno la media delle chiamate effettuate ogni mese si abbassa alle 300. Ottobre del primo anno è il mese col maggior tasso di successo (61\%), con circa 1500 chiamate effettuate.
Gli altri mesi di successo sono tutti concentrati nel secondo anno, come marzo del “terzo anno”, ma ipotizziamo che questo sia dovuto al minor numero di chiamate effettuate a quel punto.

\begin{figure}[H]
\centering
\includegraphics[width=0.6\textwidth]{GRAFICI PER LATEX/monthly-calls.png}
\caption{Distribuzione chiamate effettuate mensilmente}
\label{fig:monthly-calls}
\end{figure}

\subsubsection{Correlazione e Feature Selection}

Dallo studio dei coefficienti di correlazione, emerge effettivamente una relazione mediamente forte di segno positivo fra ‘duration’ e la variabile target. 

\begin{table}[H]
    \centering
    \begin{tabular}{|l|r|r|}
    \hline
    Metodo    & Pearson & Spearman\\
    \hline
    duration  & 0.394521 & 0.342469\\
    \hline
    pdays     & 0.103621 & 0.154055\\
    \hline
    previous  & 0.093236 & 0.169124\\
    \hline
    balance   & 0.052838 & 0.100295\\
    \hline
    age       & 0.025155 & -0.008750\\
    \hline
    day       & -0.028348 & -0.029548\\
    \hline
    campaign  & -0.073172 & -0.084054\\
         \hline
    \end{tabular}
    \caption{Coefficienti di correlazione con y}
    \label{tab:correlation-table}
\end{table}
% li metterei con 4 decimali, perché è più difficile da leggere così imo
%oi! sisi poi si può cambiare %okiii

Dallo studio dei coefficienti di correlazione, emerge effettivamente una relazione mediamente forte di segno positivo fra ‘duration’ e la variabile target. Questa relazione è la più marcata rilevata nel dataset. Possiamo ipotizzare che chiamate più lunghe equivalgano ad un dialogo prolungato sulle condizioni e le caratteristiche dell’offerta, oltre che alla finalizzazione del contratto.
È fondamentale sottolineare che questa non sia una variabile propriamente predittiva, poiché è un dato che viene registrato solo in seguito alla comunicazione col cliente e non è conoscibile in anticipo. 

Notiamo differenze per le due tipologie di indici: Pearson individua valori positivi, anche se meno forti, con la frequenza dei contatti e l’essere già stati contattati.
Sapendo che i valori dei coefficienti di correlazione variano tra -1 e 1, i valori inferiori a quelli menzionati sono trascurabili.
Spearman, d’altra parte, cattura valori leggermente più alti per quasi tutte le feature, con l’eccezione di ‘duration’ (-0,052) e di ‘age’ (-0,033).
In questo caso il numero di contatti nella campagna precedente per singolo cliente è più influente (+0,075)
La correlazione fra saldo in banca e variabile target raddoppia, pur rimanendo piuttosto bassa (0,1).
Questi sono gli effetti della robustezza di Spearman agli outlier, evidentemente presenti nelle variabili considerate.

Dal momento che i coefficienti erano stati calcolati prima di creare la variabile 'month\_index', riportiamo i risultati ottenuti dallo studio della feature importance rilevata dal modello Random Forest. In questo caso possiamo notare che sia la durata della chiamata e sia il mese in cui è stata effettuata sono risultate le variabili più discriminanti nelle scelte degli alberi. 
Tuttavia, si tratta di fattori poco sfruttabili dal punto di vista predittivo. Da un lato la durata delle chiamate che non è conoscibile a priori; dall'altro il mese in cui la chiamata viene effettuata, la cui influenza sul valore della variabile target è interpretabile solo nel contesto della campagna marketing svolta dalla banca portoghese.

\subsection{Fase 3: Confronto tra modelli}
Utilizzando tre diversi modelli di classificazione binaria e confrontando gli indicatori di performance, osserviamo che il forte \textit{sbilanciamento} delle classi, per cui la classe 0 è nettamente più numerosa della classe 1, ha inficiato le capacità di predizione di tutti i modelli (\textit{support} per la classe 1 è di appena duemila casi, contro i quasi sedicimila della classe 0). Difatti tutti i modelli hanno performance considerevolmente peggiori per la classe 1, ovvero i potenziali clienti aderenti.


\begin{table}[H] 
    \centering
    \begin{tabular}{|r|l|l|l|l|l|}
    \hline
    Modello     & Accuracy & Precision & Recall & F1-score & Support\\
    \hline
    Regressione logistica     & 0.906276 & 0.667462 & 0.396503 & 0.497480 & 2116\\
    \hline
    K-Nearest Neighbors    & 0.891125 & 0.552613 & 0.364839 & 0.439510 & 2116\\
    \hline
    Random Forest     & 0.899641 & 0.553350 & 0.737713 & 0.632368 & 2116\\
    \hline
    \end{tabular}
    \caption{Confronto performance modelli, valori classe 1}
    \label{tab:models}
\end{table}

\begin{itemize}
    \item l’\textit{accuratezza} è sostanzialmente equivalente per Regressione Logistica e Random Forest: entrambi i modelli sono stati in grado di predire correttamente nove casi su dieci nel dataset di test (ma la maggioranza di questi casi sono clienti classificati come non iscritti ad un conto deposito)
\end{itemize}

\begin{itemize}
    \item la Regressione Logistica è il modello più \textit{preciso} nell’individuare i clienti in classe 1: identifica il 66\% dei casi positivi
\end{itemize}

\begin{itemize}
    \item per la \textit{recall} per la classe 1 il modello con risultati migliori è Random Forest: ha identificato il 73\% dei clienti effettivamente iscritti, contro il 36\% degli altri due modelli
\end{itemize}

\begin{itemize}
    \item è sempre Random Forest ad ottenere il valore di \textit{f1-score} più alto
\end{itemize}

Concludiamo che, per il nostro dataset, Random Forest è il modello più adatto a prevedere i potenziali clienti; riteniamo che il risultato ottenuto per f1-score sia un compromesso accettabile. Resta il fatto che il modello non eccelle, ma è capace di gestire un dataset fortemente sbilanciato a favore dei clienti non iscritti, dimostrando un'ottima capacità predittiva per i clienti realmente iscritti.
&cosa comporta avere una precision così bassa e una recall più alta


\section{Conclusioni}
* abbiamo svolto analisi su questo
* abbiamo trovato come rilevanti le x variabili
* i modelli predittivi hanno avuto risultati simili e non eccelsi

* key takeaways: * cosa comportano precision e recall bassa * possibile confronto con altri paper per contestualizzare i nostri risultati * possibile ottenere risultati migliori?


\section{Bibliografia}
\bibliography{biblio}
\bibliographystyle{plainnat}

\end{document}
